\documentclass[a4paper,titlepage,bibtotoc,liststotoc,final,twoside,openright,cleardoubleempty,BCOR12mm]{scrreprt} % ABGABEFASSUNG

\usepackage[utf8]{inputenc}
\usepackage[T1]{fontenc}
\usepackage[ngerman]{babel}
\usepackage{epsfig}
\usepackage{amsmath}
\usepackage{hyperref}
\usepackage{bibgerm}
\usepackage{breakurl}
\usepackage{ae,aecompl}
\usepackage{graphicx}
\usepackage{tabularx}
\usepackage{rotating}
\usepackage{lscape}
\usepackage{subfigure}

% Abstände bei Trennzeichen für Zahlen korrekt machen
\mathcode`,="013B
\mathcode`.="613A

% Tex ist kein Zeichenprogramm...
\usepackage{tikz}
\usetikzlibrary{arrows,automata,shapes,topaths}

% Vorlage - Obere Linie, Seitenzahlen
\usepackage{scrlayer-scrpage}
\pagestyle{scrheadings}
\clearscrheadfoot
\automark[section]{chapter}
\setheadsepline{.4pt}
\ihead[]{\leftmark}
\ohead[]{\rightmark}
\ofoot[\pagemark]{\pagemark}

% Index erstellen
\usepackage{makeidx}
\makeindex

% Arbeitsbeschreibung & Makros
\newcommand{\Name}{Deniz Groenhoff}
\newcommand{\Email}{deniz.groenhoff@uni-oldenburg.de}
\newcommand{\MatrNr}{5477417}
\newcommand{\Titel}{Annotationsbasierte Einstiegserleichterung in die Entwicklung von JavaFX-Anwendungen}
\newcommand{\Department}{Department für Informatik}
\newcommand{\Abteilung}{Medieninformatik und Multimedia-Systeme}
\newcommand{\ErstGutachter}{Prof. Dr. Susanne Boll}
\newcommand{\ZweitGutachter}{Dr.-Ing. Dietrich Boles}
\newcommand{\ArbeitTyp}{Bachelorarbeit}


\newcommand{\MONTH}{%
	\ifcase\the\month
	\or Januar% 1
	\or Februar% 2
	\or März% 3
	\or April% 4
	\or Mai% 5
	\or Juni% 6
	\or Juli% 7
	\or August% 8
	\or September% 9
	\or Oktober% 10
	\or November% 11
	\or Dezember% 12
	\fi}

\title{\Titel}
\author{\Name}
\date{\today}

\begin{document}

% Fuers inhaltsvz gibts andere seitenzahlen
\pagenumbering{Roman}

% Titelseite erstellen
\begin{titlepage}
	\thispagestyle{empty}
% zentriert: 0.26cm
\hoffset=0.4cm
\begin{center}
\begin{figure}
\centering
\includegraphics[width=10cm]{Abbildungen/uniol_4c}\\
\end{figure}
\vspace{2cm}
{\bf \Department}\\Abteilung für \Abteilung\\
\vspace{4cm}
{\bf \ArbeitTyp}\\
\vspace{.5cm}
{\begin{minipage}[t]{10cm}
\centering
\Large
\Titel
\end{minipage}}\\
\vspace{3cm}
\Name\\
\vspace{.5cm}
\today\\
\vspace{3cm}
1. Gutachterin: \ErstGutachter\\
2. Gutachter: \ZweitGutachter
\end{center}

\end{titlepage}

\newpage
\thispagestyle{empty}
\quad

% Eigenstaendigkeitserklaerung
\newpage
\thispagestyle{empty}
\vspace{6.8cm}
\begin{center}
	\begin{Large}
		\textbf{Erklärung}
	\end{Large}
\end{center}

\noindent
Ich erkläre an Eides statt, dass ich die Arbeit selbständig verfasst
und keine anderen als die angegebenen Quellen und Hilfsmittel benutzt und die allgemeinen Prinzipien wissenschaftlicher Arbeit und Veröffentlichungen, wie sie in den Leitlinien guter wissenschaftlicher Praxis der Carl von Ossietzky Universität Oldenburg festgelegt sind, befolgt habe.

\vspace{3cm}
\noindent
\Name\\
Matrikelnummer \MatrNr\\
Oldenburg, den \the\day. \MONTH\ \the\year{}


% Zusammenfassung erstellen
\begin{abstract}
	\section*{Zusammenfassung}
Hier kommt in der Regel eine ca. halbseitige Zusammenfassung von Motivation und
Ergebnis der Arbeit hin. Eine zeitliche Abfolge, wann was gemacht wurde, spielt hier
keine Rolle
\footnote{Fussnote 1}
	\begin{center}
	\Large\bfseries{Abstract}
\end{center}
Abstract (jeweils auf Englisch und Deutsch)
\end{abstract}

% Inhaltsverzeichnis erstellen
\tableofcontents
% Abbildungsverzeichnis
\listoffigures
% Tabellenverzeichnis
\listoftables

\chapter{Einleitung}
\pagenumbering{arabic} % Seitenzahlen umstellen
\label{einleitung}
Einleitung

\chapter{Inhalt 1}

\chapter{Inhalt 2}

\chapter{Fazit}

\begin{appendix}
\chapter{Appendix 1}
\chapter{Appendix 2}
\end{appendix}

% Literaturverzeichnis...
\nocite{*}
\bibliographystyle{geralpha}
\bibliography{Bibliographie/Bibliographie}

% Index ausgeben
\newpage
\printindex

\end{document}
